\begin{tomtat}
	\begin{minipage}[t]{0.3\textwidth}
		\vspace{0pt}
		\begin{tomtat}
			\begin{center}
				\textbf{ĐỀ 50}
			\end{center}
			
		\end{tomtat}
	\end{minipage}
	\hfill
	\begin{minipage}[t]{0.5\textwidth}
		\vspace{0pt}
		\begin{tomtat}
			\begin{center}
				\includegraphics[scale=0.08]{images/50}
			\end{center}
		\end{tomtat}
	\end{minipage}
\end{tomtat}
\textbf{I. PHẦN TRẮC NGHIỆM}
\begin{ex} % Câu 1
	Giá trị của biểu thức $4\sqrt{\frac{25}{16}}+5\sqrt{\frac{9}{25}}$ là:
	\choice{$9$}{$8$}{$7$}{$\frac{5}{4}$}
	\loigiai{
		\textbf{Đáp án đúng là B.} \\
		\begin{align*}
			4\sqrt{\frac{25}{16}}+5\sqrt{\frac{9}{25}} &= 4 \cdot \frac{\sqrt{25}}{\sqrt{16}} + 5 \cdot \frac{\sqrt{9}}{\sqrt{25}} \\
			&= 4 \cdot \frac{5}{4} + 5 \cdot \frac{3}{5} \\
			&= 5 + 3 \\
			&= 8
		\end{align*}
	}
\end{ex}

\begin{ex} % Câu 2
	Nghiệm của phương trình $4x+1=5x$ là:
	\choice{$2$}{$3$}{$4$}{$1$}
	\loigiai{
		\textbf{Đáp án đúng là D.} \\
		\begin{align*}
			4x+1 &= 5x \\
			1 &= 5x - 4x \\
			1 &= x \\
			x &= 1
		\end{align*}
	}
\end{ex}

\begin{ex} % Câu 3
	Bất phương trình $5x-30<0$ có nghiệm là:
	\choice{$x<6$}{$x>6$}{$x \le 6$}{$x \ge 6$}
	\loigiai{
		\textbf{Đáp án đúng là A.} \\
		\begin{align*}
			5x-30 &< 0 \\
			5x &< 30 \\
			x &< \frac{30}{5} \\
			x &< 6
		\end{align*}
	}
\end{ex}

\begin{ex} % Câu 4
	Điểm nào sau đây thuộc đồ thị hàm số $y=3x^2$?
	\choice{$(1;3)$}{$(0;3)$}{$(-1;-3)$}{$(3;0)$}
	\loigiai{
		\textbf{Đáp án đúng là A.} \\
		Để kiểm tra một điểm có thuộc đồ thị hàm số hay không, ta thay tọa độ của điểm đó vào phương trình hàm số $y=3x^2$.
		A. Với $(1;3)$: Thay $x=1, y=3 \Rightarrow 3 = 3(1)^2 \Rightarrow 3 = 3(1) \Rightarrow 3 = 3$ (Đúng).
		B. Với $(0;3)$: Thay $x=0, y=3 \Rightarrow 3 = 3(0)^2 \Rightarrow 3 = 0$ (Sai).
		C. Với $(-1;-3)$: Thay $x=-1, y=-3 \Rightarrow -3 = 3(-1)^2 \Rightarrow -3 = 3(1) \Rightarrow -3 = 3$ (Sai).
		D. Với $(3;0)$: Thay $x=3, y=0 \Rightarrow 0 = 3(3)^2 \Rightarrow 0 = 3(9) \Rightarrow 0 = 27$ (Sai).
		Vậy, điểm $(1;3)$ thuộc đồ thị hàm số.
	}
\end{ex}

\begin{ex} % Câu 5
	Gọi $l, h, r$ lần lượt là độ dài đường sinh, chiều cao và bán kính của hình trụ (T). Thể tích $V$ của hình trụ (T) là:
	\choice{$V=\frac{4}{3}\pi r^3$}{$V=\frac{1}{3}\pi r^2l$}{$V=\pi r^2h$}{$V=\frac{4}{3}\pi r^2h$}
	\loigiai{
		\textbf{Đáp án đúng là C.} \\
		Công thức tính thể tích của hình trụ là $V = \pi r^2h$, trong đó:
		\begin{itemize}
			\item $r$ là bán kính đáy của hình trụ.
			\item $h$ là chiều cao của hình trụ.
		\end{itemize}
		Các lựa chọn khác:
		A. $V=\frac{4}{3}\pi r^3$ là công thức thể tích hình cầu.
		B. $V=\frac{1}{3}\pi r^2l$ không phải là công thức chuẩn, $V=\frac{1}{3}\pi r^2h$ là công thức thể tích hình nón (với $h$ là chiều cao, $l$ là đường sinh).
		D. $V=\frac{4}{3}\pi r^2h$ không phải là công thức chuẩn.
	}
\end{ex}

\begin{ex} % Câu 6
	Tam giác ABC vuông tại A có AB = 3cm; BC = 5cm. Giá trị của $\tan B$ là:
	\choice{$\frac{3}{4}$}{$\frac{3}{5}$}{$\frac{4}{3}$}{$\frac{4}{5}$}
	\loigiai{
		\textbf{Đáp án đúng là C.} \\
		Trong tam giác ABC vuông tại A, áp dụng định lý Pitago, ta có:
		$AC^2 = BC^2 - AB^2$
		\begin{align*}
			AC^2 &= 5^2 - 3^2 \\
			AC^2 &= 25 - 9 \\
			AC^2 &= 16 \\
			AC &= \sqrt{16} = 4 \text{ (cm)}
		\end{align*}
		Trong tam giác ABC vuông tại A, ta có:
		$\tan B = \frac{\text{cạnh đối}}{\text{cạnh kề}} = \frac{AC}{AB}$
		\begin{align*}
			\tan B &= \frac{4}{3}
		\end{align*}
	}
\end{ex}
\begin{ex} % Câu 7
	Tuổi nghề (đơn vị: năm) của 33 giáo viên ở một trường trung học cơ sở được ghi lại như sau:
\begin{center}
	\begin{tabular}{|c|c|c|c|c|c|c|}
		\hline
		7 & 7 & 10 & 3 & 4 & 4 & 9 \\
		\hline
		3 & 3 & 7 & 7 & 12 & 4 & 3 \\
		\hline
		12 & 9 & 12 & 9 & 10 & 7 & 7 \\
		\hline
		12 & 7 & 10 & 4 & 10 & 10 & \\
		\hline
		4 & 9 & 7 & 10 & 9 & 8 & \\
		\hline
	\end{tabular}
\end{center}
	Tần số của giá trị $x=7$ là:
	\choice{$n=7$}{$n=8$}{$n=6$}{$n=5$}
	\loigiai{
		\textbf{Đáp án đúng là B.} \\
		Để tìm tần số của giá trị $x=7$, ta đếm số lần xuất hiện của số 7 trong bảng dữ liệu:
		\begin{itemize}
			\item Hàng 1: có 4 số 7 (vị trí 1, 2, 10, 11).
			\item Hàng 2: có 2 số 7 (vị trí 9, 10).
			\item Hàng 3: có 2 số 7 (vị trí 1, 8).
		\end{itemize}
		Tổng số lần xuất hiện của giá trị 7 là $4 + 2 + 2 = 8$.
		Vậy, tần số của giá trị $x=7$ là $n=8$.
	}
\end{ex}

\begin{ex} % Câu 8
	Gieo một con xúc xắc cân đối và đồng chất. Gọi biến cố A: “số chấm xuất hiện trên con xúc xắc là số lẻ”. Số kết quả thuận lợi cho biến cố A là:
	\choice{$6$}{$2$}{$4$}{$3$}
	\loigiai{
		\textbf{Đáp án đúng là D.} \\
		Khi gieo một con xúc xắc cân đối và đồng chất, các mặt có thể xuất hiện là $\{1, 2, 3, 4, 5, 6\}$.
		Biến cố A: “số chấm xuất hiện trên con xúc xắc là số lẻ”.
		Các số lẻ trong các mặt của xúc xắc là: $1, 3, 5$.
		Vậy, có 3 kết quả thuận lợi cho biến cố A.
	}
\end{ex}

\newpage
\textbf{II. PHẦN TỰ LUẬN}
\begin{ex} % Câu 9
	Rút gọn biểu thức $A=\left(\frac{1}{x-4}+\frac{1}{x+4\sqrt{x}+4}\right) \cdot \frac{x+2\sqrt{x}}{\sqrt{x}}$ (với $x>0, x \ne 4$)
	\loigiai{
		Điều kiện: $x>0, x \ne 4$.
		Ta có:
		$x-4 = (\sqrt{x}-2)(\sqrt{x}+2)$
		$x+4\sqrt{x}+4 = (\sqrt{x}+2)^2$
		Xét phần trong ngoặc:
		\begin{align*}
			\frac{1}{x-4}+\frac{1}{x+4\sqrt{x}+4} &= \frac{1}{(\sqrt{x}-2)(\sqrt{x}+2)} + \frac{1}{(\sqrt{x}+2)^2} \\
			&= \frac{\sqrt{x}+2 + (\sqrt{x}-2)}{(\sqrt{x}-2)(\sqrt{x}+2)^2} \\
			&= \frac{\sqrt{x}+2+\sqrt{x}-2}{(\sqrt{x}-2)(\sqrt{x}+2)^2} \\
			&= \frac{2\sqrt{x}}{(\sqrt{x}-2)(\sqrt{x}+2)^2}
		\end{align*}
		Biểu thức $\frac{x+2\sqrt{x}}{\sqrt{x}}$:
		$$ \frac{x+2\sqrt{x}}{\sqrt{x}} = \frac{\sqrt{x}(\sqrt{x}+2)}{\sqrt{x}} = \sqrt{x}+2 $$
		(Vì $x>0$ nên $\sqrt{x} \ne 0$).
		Vậy biểu thức $A$ trở thành:
		\begin{align*}
			A &= \frac{2\sqrt{x}}{(\sqrt{x}-2)(\sqrt{x}+2)^2} \cdot (\sqrt{x}+2)
		\end{align*}
		Vì $x \ne 4$ nên $\sqrt{x}+2 \ne 0$. Ta có thể rút gọn:
		\begin{align*}
			A &= \frac{2\sqrt{x}}{(\sqrt{x}-2)(\sqrt{x}+2)} \\
			&= \frac{2\sqrt{x}}{x-4}
		\end{align*}
		Vậy, với $x>0, x \ne 4$ thì $A = \frac{2\sqrt{x}}{x-4}$.
	}
\end{ex}

\begin{ex} % Câu 10
	Biết phương trình bậc hai $x^2+4x-3=0$ có hai nghiệm phân biệt $x_1, x_2$. Không giải phương trình, tính giá trị của biểu thức: $A=x_1^2x_2+x_1x_2^2-x_1-x_2$.
	\loigiai{
		Phương trình đã cho là $x^2+4x-3=0$.
		Theo hệ thức Vi-ét, ta có:
		$ \begin{cases} x_1 + x_2 = - \frac{4}{1} = -4 \\ x_1 x_2 = \frac{-3}{1} = -3 \end{cases} $
		Ta biến đổi biểu thức A:
		\begin{align*}
			A &= x_1^2x_2+x_1x_2^2-x_1-x_2 \\
			&= x_1x_2(x_1+x_2) - (x_1+x_2) \\
			&= (x_1x_2-1)(x_1+x_2)
		\end{align*}
		Thay giá trị từ hệ thức Vi-ét vào biểu thức A:
		\begin{align*}
			A &= (-3-1)(-4) \\
			&= (-4)(-4) \\
			&= 16
		\end{align*}
		Vậy, giá trị của biểu thức $A = 16$.
	}
\end{ex}

\begin{ex} % Câu 11
	Giải hệ phương trình: $ \begin{cases} 2x-3y=-5 \\ x+3y=11 \end{cases} $
	\loigiai{
		Ta có hệ phương trình:
		$ \begin{cases} 2x-3y=-5 \quad &(1) \\ x+3y=11 \quad &(2) \end{cases} $
		Cộng vế theo vế của phương trình (1) và (2), ta được:
		\begin{align*}
			(2x-3y) + (x+3y) &= -5+11 \\
			2x - 3y + x + 3y &= 6 \\
			3x &= 6 \\
			x &= 2
		\end{align*}
		Thay $x=2$ vào phương trình (2):
		\begin{align*}
			2 + 3y &= 11 \\
			3y &= 11 - 2 \\
			3y &= 9 \\
			y &= 3
		\end{align*}
		Vậy, nghiệm của hệ phương trình là $(x;y) = (2;3)$.
	}
\end{ex}
\begin{ex} % Câu 12a - Xác suất
	Viết ngẫu nhiên một số tự nhiên có hai chữ số lớn hơn 50. Xác suất của biến cố A: “Số tự nhiên được viết là bội của 6” là bao nhiêu?
	\loigiai{
		Các số tự nhiên có hai chữ số lớn hơn 50 là các số từ 51 đến 99.
		Số lượng các số này là $99 - 51 + 1 = 49$ số.
		Vậy, số phần tử của không gian mẫu là $n(\Omega) = 49$.
		Vậy, các số thuận lợi cho biến cố A là $\{54, 60, 66, 72, 78, 84, 90, 96\}$.
		Số phần tử của biến cố A là $n(A) = 8$.
		
		Xác suất của biến cố A là:
		$P(A) = \frac{n(A)}{n(\Omega)} = \frac{8}{49}$.
	}
\end{ex}

\begin{ex} % Câu 12b - Bài toán ủng hộ sách
	Nhân dịp Ngày sách và Văn hóa đọc Việt Nam, hai lớp 9A và 9B cùng tham gia ủng hộ sách cho thư viện. Tổng số học sinh của hai lớp là 82. Mỗi học sinh lớp 9A ủng hộ 3 quyển sách, còn mỗi học sinh lớp 9B ủng hộ 4 quyển sách. Tổng số sách hai lớp ủng hộ là 288 quyển. Hỏi mỗi lớp có bao nhiêu học sinh?
	\loigiai{
		Gọi $x$ là số học sinh của lớp 9A.
		Gọi $y$ là số học sinh của lớp 9B.
		(Điều kiện: $x, y$ là các số nguyên dương)
		
		Tổng số học sinh của hai lớp là 82, nên ta có phương trình:
		$x + y = 82 \quad (1)$
		
		Số sách lớp 9A ủng hộ là: $3x$ (quyển).
		Số sách lớp 9B ủng hộ là: $4y$ (quyển).
		Tổng số sách hai lớp ủng hộ là 288 quyển, nên ta có phương trình:
		$3x + 4y = 288 \quad (2)$
		
		Từ (1) và (2), ta có hệ phương trình:
		$ \begin{cases} x + y = 82 \\ 3x + 4y = 288 \end{cases} $
		Từ phương trình (1), suy ra $y = 82 - x$.
		Thế $y = 82 - x$ vào phương trình (2):
		\begin{align*}
			3x + 4(82 - x) &= 288 \\
			3x + 328 - 4x &= 288 \\
			-x &= 288 - 328 \\
			-x &= -40 \\
			x &= 40
		\end{align*}
		Thay $x = 40$ vào $y = 82 - x$:
		\begin{align*}
			y &= 82 - 40 \\
			y &= 42
		\end{align*}
		Vậy, lớp 9A có 40 học sinh và lớp 9B có 42 học sinh.
	}
\end{ex}
\begin{ex} % Câu 13
	Cho $\triangle ABC$ nhọn có $AB > AC$, nội tiếp đường tròn $(O)$, các đường cao $AD, CE$ cắt nhau tại $H$ ($D \in BC, E \in AB$). Qua $O$ kẻ đường thẳng song song với $AD$, cắt $BC, AB, AC$ lần lượt tại $M, N$ và $P$.
	\begin{enumerate}[a)]
		\item Chứng minh rằng tứ giác $ACDE$ nội tiếp đường tròn và $\widehat{APN} = \widehat{CED}$.
		\item Gọi $I$ là trung điểm của $PN$. Chứng minh rằng $AP \cdot BM = PI \cdot HB$ và $IA$ vuông góc với $MH$.
	\end{enumerate}
	\loigiai{}
\end{ex}

\begin{ex} % Câu 14
	\begin{enumerate}[a)]
		\item Một công ty sản xuất dụng cụ thể thao nhận được một đơn đặt hàng sản xuất 8000 quả bóng tennis. Công ty này sở hữu một số máy móc, mỗi máy móc có thể sản xuất 30 quả bóng trong một giờ. Chi phí thiết lập các máy này là 200000 đồng cho mỗi máy. Khi được thiết lập, hoạt động sản xuất sẽ hoàn toàn diễn ra tự động dưới sự giám sát. Số tiền phải trả cho người giám sát là 192000 đồng một giờ (người này sẽ giám sát tất cả các máy hoạt động). Số máy móc công ty nên sử dụng là bao nhiêu để chi phí sản xuất là thấp nhất?
		\item Cho $x, y, z$ là độ dài 3 cạnh của một tam giác. Chứng minh rằng:
		$\sqrt{x+y-z} + \sqrt{y+z-x} + \sqrt{z+x-y} \le \sqrt{x} + \sqrt{y} + \sqrt{z}$.
	\end{enumerate}
	\loigiai{}
\end{ex}

\setcounterref{ex}{0}
